\documentclass[svgnames]{article}     % use "amsart" instead of "article" for AMSLaTeX format
%\geometry{landscape}                 % Activate for rotated page geometry

%\usepackage[parfill]{parskip}        % Activate to begin paragraphs with an empty line rather than an indent

\usepackage{graphicx}                 % Use pdf, png, jpg, or eps§ with pdflatex; use eps in DVI mode

%maths                                % TeX will automatically convert eps --> pdf in pdflatex
\usepackage{amssymb}
\usepackage{amsmath}
\usepackage{esint}
\usepackage{geometry}

% Inverting Color of PDF
\usepackage{xcolor}
\pagecolor[rgb]{0.19,0.19,0.19}
\color[rgb]{0.77,0.77,0.77}

%noindent
\setlength\parindent{0pt}

%pgfplots
\usepackage{pgfplots}

%images
%\graphicspath{{ }}                   % Activate to set a image directory

%tikz
\usepackage{pgfplots}
\pgfplotsset{compat=1.15}
\usepackage{comment}
\usetikzlibrary{arrows}
\usepackage[most]{tcolorbox}

\usepackage{listings}
\usepackage{color}

\definecolor{dkgreen}{rgb}{0,0.6,0}
\definecolor{gray}{rgb}{0.5,0.5,0.5}
\definecolor{mauve}{rgb}{0.58,0,0.82}

\lstset{frame=tb,
  language=Java,
  aboveskip=3mm,
  belowskip=3mm,
  showstringspaces=false,
  columns=flexible,
  basicstyle={\small\ttfamily},
  numbers=none,
  numberstyle=\tiny\color{gray},
  keywordstyle=\color{blue},
  commentstyle=\color{dkgreen},
  stringstyle=\color{mauve},
  breaklines=true,
  breakatwhitespace=true,
  tabsize=3
}

\lstset{language=Python}

\tcbset {
  base/.style={
    arc=0mm,
    bottomtitle=0.5mm,
    boxrule=0mm,
    colbacktitle=black!10!white,
    coltitle=black,
    fonttitle=\bfseries,
    left=2.5mm,
    leftrule=1mm,
    right=3.5mm,
    title={#1},
    toptitle=0.75mm,
  }
}

\definecolor{brandblue}{rgb}{0.34, 0.7, 1}
\newtcolorbox{mainbox}[1]{
  colframe=brandblue,
  base={#1}
}

\newtcolorbox{subbox}[1]{
  colframe=black!30!white,
  base={#1}
}

%Figures
\usepackage{float}
\usepackage{caption}
\usepackage{lipsum}


\title{Brief Article}
\author{Deval Deliwala}
%\date{}                              % Activate to display a given date or no date

\begin{document}
\maketitle
%\section{}
%\subsection{}
\tableofcontents                     % Activate to display a table of contents
\newpage

We can always choose to view multiple systems \textit{together} as if they form
a single, compound system -- to which the discussion in the previous section
applies. 

\section{Classical Information}

Recognizing how mathematics works in the familiar setting of classical
information is helpful in understanding why quantum information is described
the way it is. 

\subsection{Classical States via the Cartesian Product}

We will start with classical states of multiple systems. We will begin by
discussing just two systems and then generalize further. \\

To be precise, let us suppose that $X$ is a system whose classical state set is
$\Sigma$ and $Y$ is a second system having classical state set $\Gamma$.
Because these are \textit{classical state sets}, our assumption is that
$\Sigma$ and $\Gamma$ are both finite and nonempty. It could be that $\Sigma
= \Gamma$, but this is not necessarily so. \\

Now imagine the two systems  $X$ and $Y$ are placed side-by-side, with $X$ on the left and $Y$ on the right. We can views these two systems as if they form a single system $(X, Y)$ or $XY$. \\

The set of classical states of $XY$ is the \textit{Cartesian Product} of $\Sigma$ and $\Gamma$, which is the set defined as 

\begin{mainbox}{Cartesian Product}
\[
  \Sigma \times \Gamma = { (a, b) : a \in \Sigma \text{ and } b \in \Gamma}
\] 
\end{mainbox}

The Cartesian Product is the mathematical notion of viewing an element of one set and an element of a second set together, as if they form a single set. \\ 

For more than two systems, the situation generalizes in a natural way. If we suppose that $X_1, \cdots, X_n$ are systems having classical state sets $\Sigma_1, \cdots, \Sigma_n$, respectively for any positive integer $n$, the classical state set of the $n$-tuple $X_1X_2\cdots X_n$, viewed as a single joint system, is the Cartesian Product

\[
  \Sigma_1 \times \cdots \times \Sigma_n = {(a_1, \cdots, a_n) : a_1 \in \Sigma_1 \cdots a_n \in \Sigma_n}
\] \vspace{5px}

\subsection{Representing states as strings} 

It is convenient to write a classical state $(a_1, \cdots, a_n)$ as a string `$a_1 \cdots a_n$`. The notion of a string, which is fundamentally important concept in CS, is formalized in mathematical terms through the cartesian product. \\

For example suppose that $X_1, \cdots , X_{10}$ are bits, so the classical state sets of these systems are all the same 

\[
  \Sigma_1 = \Sigma_2 = \cdots = \Sigma_{10} = {0, 1}
\] \vspace{5px}

The set ${0, 1}$ is referred to as the \textit{binary alphabet}. There are
$2^{10} = 1024$ classical states of the joint system $X_1X_2 \cdots X_{10}$
which are elements of the set

\[
  \Sigma_1 \times \Sigma_2 \times \cdots \times \Sigma_{10} = {0, 1}^10. 
\] \vspace{5px}

Written as strings, these classical states look like this: 

\begin{align*}
  00000&00000 \\
  00000&00001 \\
  00000&00010 \\
  00000&00011 \\
  00000&00100 \\
       &\vdots \\
  11111&11111
\end{align*}

For the classical state $0001010000$ we see that $X_4$ and $X_6$ are in the
state 1, while all other systems are in the state 0. 





\end{document}

